% Layout extracted from Kupzyk report at ca1d062; retain a descriptive title.
\documentclass[11pt]{article}
\usepackage[letterpaper,margin=0.82in]{geometry}
\usepackage{booktabs,longtable,array,graphicx,amsmath,xcolor,hyperref,xurl}
\usepackage[authoryear,round]{natbib}
\hypersetup{colorlinks=true,urlcolor=blue,linkcolor=blue,citecolor=blue}
\renewcommand{\arraystretch}{1.25}
\setlength{\emergencystretch}{1em}
\setlength{\parskip}{5pt}\setlength{\parindent}{0pt}
% Use scalable fonts; inspect pdffonts output after building. Avoid adding T1
% fontenc without an installed scalable font family. Add title/date in manuscript.

\hypersetup{pdftitle={Mediation in Regression Discontinuity Designs: A Replication of Chen (2021)},pdfauthor={vibe-math replication project}}
\usepackage{enumitem}
\usepackage{needspace}
\def\UrlBreaks{\do\/\do\_\do\-}\def\UrlBigBreaks{}
\setlist{itemsep=1pt,topsep=2pt}
\newcommand{\cp}{c$'$}

\title{\vspace*{-3.6em}Mediation in Regression Discontinuity Designs:\\ A Replication of Chen (2021)}
\author{}
\date{}

\begin{document}
\maketitle
\vspace{-5em}

\begin{abstract}
\noindent We did not replicate three of the four simulation headlines of \citet{chen2021}, who reported biased direct-effect estimates, percentile bootstrap intervals that cover the true value far too rarely, and power for the mediated interaction (MINT) near zero. Of her direct-effect results, 85.42\% fall outside [-.10, .10] against 0.00\% of ours, and her mean coverage is 43.62 points below 95\% against 0.04. At n = 500 for the largest interaction, at her two directly read cutoffs, her MINT rejection rate exceeds the nominal .05 by 0.013 (raw 0.063) against 0.802 (raw 0.852). Her tentative minimum of 5,000 cases for accurate Type I error and power above .80 holds in part: at n = 5,000, 86.3\% of zero-effect results are in Bradley's Type I range and 89.2\% of nonzero results reach .80 power, but products with both paths zero stay conservative. Her Head Start estimates need restricted-use data and could not be rerun; from her printed numbers, her arithmetic recomputes and her conclusion of opposite significance holds in 10 of 15 comparisons.\par\smallskip
\noindent Her bootstrap is unseeded, so interval results use our seeds. Her printed, seeded R code, replayed on the 1,280,000 datasets it generates, recovers the interaction and is unbiased like our run. Her figures instead match results from data without the treatment-by-mediator interaction. Her coverage matches a comparison with population values wrong for treated cases, while the rule behind her bias is not determined. Of ten disagreements, eight are explained and two partly. Table 1 gives the tolerances (her printed rounding or our figure-reading error); comparisons are of point estimates, not equivalence.
\end{abstract}

\vspace{2pt}
\noindent{\small\textbf{Table 1.} Headline comparison with \citet{chen2021} (Tables 4.8-4.9, Figure 4.28, Chapter 5).}\par\vspace{2pt}
{\scriptsize
\setlength{\tabcolsep}{3pt}\renewcommand{\arraystretch}{1.0}%
\hyphenpenalty=10000 \exhyphenpenalty=10000
\noindent\begin{tabular}{@{}>{\raggedright\arraybackslash}p{52pt}>{\raggedright\arraybackslash}p{98pt}>{\raggedright\arraybackslash}p{106pt}>{\raggedright\arraybackslash}p{41pt}>{\raggedright\arraybackslash}p{41pt}>{\raggedright\arraybackslash}p{42pt}>{\raggedright\arraybackslash}p{41pt}>{\raggedright\arraybackslash}p{31pt}@{}}
\toprule
Effect and estimand & Published & Ours & Published effect & Our effect & Effect magnitude & Direction & Verdict\\
\midrule
E1 (H1). Direct effects outside [-.10,~.10] & PNDE and TNDE within [-.10, .10] in only 12.50-18.75\% of conditions, with negative bias & PNDE and TNDE within [-.05, .05] in 100\% of conditions; largest $|$bias$|$ 0.016 & 85.42 & 0.00 & smaller & zero vs nonzero & no match\\
\midrule
E2 (H2). Coverage short of 95\% & Coverage 24.10-83.99\% by effect and cutoff; true value usually above the CI & Coverage 94.6-95.6\% by effect and cutoff; misses on both sides at similar rates & 43.62 & 0.04 & smaller & same & no match\\
\midrule
E3 (H2). Above minus below the CI & True value above the CI in 14.08-75.55\% of replications and below it in 0.34-1.93\%, by effect and cutoff & True value above the CI in 2.5-2.9\% of replications and below it in 1.8-2.7\%, by effect and cutoff & 47.07 & 0.33 & smaller & same & no match\\
\midrule
E4 (H3). MINT power minus .05, n~=~500, ah~=~.20, cutoffs 2-3 & MINT power near zero at every sample size; never reaches .80 (Figure 4.28: 0.057-0.068 in this cell) & MINT power 0.852 at n = 500 for ah = .20 at cutoffs 2-3 (0.836 over all cutoffs); at least .97 for ah $>$= .13 from n = 1,000 & 0.013 & 0.802 & larger & same & no match\\
\midrule
E5 (H4). Minimum n; qualitative, no scalar contrast & Tentative minimum n of 5,000 for accurate Type I error and power above .80 & At n = 5,000, 86.3\% of zero-effect results have Type I in range and 89.2\% of nonzero results reach .80 power, similar at each cutoff; both-zero products stay conservative & not established & not established & not established & not established & partial\\
\bottomrule
\end{tabular}\par\vspace{2pt}
\noindent\parbox{\textwidth}{\raggedright Descriptive point estimates; not equivalence. E1: share of PNDE and TNDE condition-cutoff results with $|$bias$|$ $>$ .10 (six Table 4.8 rows); E2: 95 minus mean coverage; E3: share with the true value above minus below the interval (E2-E3 over 15 Table 4.9 rows); E1-E3 in percentage points. E4: rejection rate over b, \cp{} and cutoffs 2-3 minus the nominal .05 (raw 0.063 and 0.852; her cutoff-1 series is hidden; cell chosen after the results were seen). Positive: more bias, less coverage, value more often above, zero excluded more often. E1 tolerance 0.005 percentage points; E2 tolerance 0.005 percentage points; E3 tolerance 0.01 percentage points (rounding); E4 tolerance 0.0091 rejection rate (mean reading error of her two read points); E5 tolerance not applicable. Magnitude: absolute sizes beyond the tolerance. Monte Carlo errors: Section~\ref{sec:intro}; Chen reports none. No match = not replicated; partial = partly replicated.}
}

\setcounter{table}{1}

\section{Introduction}\label{sec:intro}

We did not replicate three of the four simulation headlines of \citet{chen2021}: her finding that the natural direct effects are biased (85.42\% of direct-effect results outside [-.10, .10] in her Table 4.8, 0.00\% in ours), her finding that 95\% percentile bootstrap intervals cover the population value too rarely with the value usually above the interval (coverage 24.10-83.99\% by effect and cutoff in her Table 4.9, 94.6-95.6\% in ours), and her finding that power for the mediated interaction stays near zero (at n = 500 for the largest interaction and her two directly read cutoffs, a rejection rate of 0.063 in her figure against 0.852 in ours, that is 0.013 against 0.802 above the nominal .05). A rate near .05 or .95 from 1,000 replications has a Monte Carlo standard error of about .007, and our E4 value one of about .002; Chen reports none. Her fourth headline, the tentative recommendation of at least 5,000 cases, holds for most but not all cells of our run (86.3\% of zero-effect results with Type I error in Bradley's range and 89.2\% of nonzero results with power of at least .80 at n = 5,000). The limit of our evidence is that her Head Start analysis (research question 2) could not be rerun without restricted-use data (from her printed numbers, her arithmetic recomputes and her conclusion of opposite significance is partly replicated, holding in 10 of 15 comparisons; S9, one of the supporting claims S1-S9 listed with H1-H4 in Appendix~\ref{app:inventory}), and that her interval-based results cannot be reproduced draw for draw, because her bootstrap sets no seed. No estimate reverses its sign. Our direct-effect results are all inside the band where 85.42\% of hers are outside; her MINT rejection rate exceeds the nominal .05 by 0.013, about 1.4 reading errors, where ours exceeds it by 0.802. Our coverage shortfall (0.04 points against her 43.62) is about three orders of magnitude smaller than hers and our imbalance (0.33 against 47.07) about two; as point estimates they point the same way as hers, with Monte Carlo standard errors of about 0.005 points if cells are independent and 0.019 if they are correlated within a condition (\href{https://github.com/alectionary/vibe-math/blob/master/formalization/chen-rdd-mediation/results/mc_se_headline.json}{\nolinkurl{results/mc_se_headline.json}}, code \href{https://github.com/alectionary/vibe-math/blob/master/formalization/chen-rdd-mediation/secondary/mc_se_headline.py}{\nolinkurl{secondary/mc_se_headline.py}}).

Tolerances are stated with Table~1: half a unit in her last printed decimal for her tables, and the reading error of her figure for the MINT power value. These magnitude labels compare point estimates; they are not tests of equivalence. Monte Carlo error is reported with each result.

\citet{chen2021} asked three research questions (her Chapter 1): which assumptions mediation analysis in a regression discontinuity design (RDD) needs (RQ1, a literature synthesis in Table 4.1); whether mediation estimates from a synthetic RDD built from the Head Start Impact Study (HSIS) agree with those from the randomized trial (RQ2, Tables 4.2-4.7); and how mediation estimates behave in a simulated sharp RDD across sample size, path effect sizes and cutoff (RQ3, Tables 4.8-4.9 and Figures 4.12-4.38). Her Chapter 5 draws four conclusions from RQ3, which are the headline claims H1-H4 of this replication. We ran her full RQ3 design, replayed her printed, seeded R code on all 1,280,000 datasets that code generates, digitized her result figures, recomputed every printed number we could from her Head Start tables, and tested competing explanations for each disagreement.

Our replay of her printed code behaves like our independent run, not like her published results. The disagreements therefore arise in steps or runs she did not print. Two patterns explain most of them: her figures match results computed from data in which the treatment-by-mediator interaction $h$ was absent (Section~\ref{dis:h-absent}), and her coverage matches a comparison with population values that are right for control cases and wrong for treated cases, while the rule behind her bias is not determined (Section~\ref{dis:population-values}, only partly explained).

\section{What we did, and what we took from others}\label{sec:contrib}

The model, the simulation design, the effect definitions and every published number are Chen's \citep{chen2021}. Her data-generating and estimation code is printed in her Appendices A-B (printed pp. 175-189). Her effect definitions follow the potential-outcomes framework for mediation with a treatment-by-mediator interaction \citep{valeri2013,vanderweele2015}. Her robustness criteria are Bradley's intervals for Type I error and coverage \citep{bradley1978} and the bias band of \citet{hoogland1998}.

Our contributions are these. (1) An independent Python/NumPy re-implementation of her Appendices A-B, run over the full factorial of her Table 3.1 at her printed 1,000 replications and 500 bootstrap resamples (our simulation). (2) An exact replay of her own printed, seeded R code on all 1,280,000 datasets it generates (our calculation with her code). (3) Digitized values of her Figures 4.12-4.38 (our reading of her figures). (4) A cell-by-cell check of 46 testable Type I error and power statements in her Chapter 4 (our reading of her text, with numeric thresholds we supplied). (5) Recomputation of her printed Head Start arithmetic, identities and significance conclusions (our calculation from her printed numbers). (6) Tests of competing explanations for ten disagreements (our analysis, recorded in \href{https://github.com/alectionary/vibe-math/blob/master/formalization/chen-rdd-mediation/DISAGREEMENTS.yaml}{\nolinkurl{DISAGREEMENTS.yaml}}). Throughout, a number is labeled as hers (with a page, table or figure locator), as our simulation or calculation, or as an assumption we supplied (Section~\ref{sec:assumptions}).

\section{Methods}\label{sec:methods}

\subsection{Design and data-generating model}\label{sec:design}

The question here is what was simulated; the answer is Chen's sharp RDD with one mediator, exactly as her Appendix A codes it. For each replication, an assignment variable $x$ is drawn from $N(0,1)$. At each of three cutoffs, the 25th, 50th and 75th sample percentile of $x$, treatment is $t = 1$ for $x$ at or above the cutoff, the centered assignment variable is $r = x - \text{cutoff}$, and new residuals are drawn:
\[ m = a\,t + e_m,\quad e_m \sim N(0, 1-a^2); \qquad y = c'\,t + b\,m + h\,t\,m + e_y,\quad e_y \sim N(0, 1-b^2), \]
where the second argument of $N$ is the variance. The fitted models are $m \sim 1 + t + r$ and $y \sim 1 + t + r + m + tm$ by ordinary least squares. The factors (her Table 3.1) are $n \in \{500, 1000, 2500, 5000, 10000\}$, $a, b, c' \in \{0, .14, .39, .59\}$ and $h \in \{0, .05, .15, .34\}$, giving 1,280 data conditions, each analysed at three cutoffs: 3,840 condition-cutoff cells, all of which were run. Cutoff index 1 in her code is the 25th percentile, so 75\% of cases are treated; we report cutoffs by the code (Section~\ref{sec:amb-cutoff}).

\subsection{Effects, population values and outcome measures}\label{sec:effects}

Five effects are estimated from the fitted coefficients: the pure natural indirect effect PNIE $= \hat a\hat b$, the total natural indirect effect TNIE $= \hat a\hat b + \hat a\hat h$, the pure natural direct effect PNDE $= \hat c' + \hat h\,\hat\imath_{0M}$, the total natural direct effect TNDE $= \hat c' + \hat h\,\hat\imath_{0M} + \hat a\hat h$, and the mediated interaction MINT $= \hat a\hat h$, where $\hat\imath_{0M}$ is the intercept of the mediator model \citep{valeri2013,vanderweele2015}. Under her data-generating model the mediator intercept is 0, so the population values are PNIE $= ab$, TNIE $= ab + ah$, PNDE $= c'$, TNDE $= c' + ah$ and MINT $= ah$ (function \texttt{population} in \href{https://github.com/alectionary/vibe-math/blob/master/formalization/chen-rdd-mediation/simulation/chen_sim.py#L41}{\nolinkurl{simulation/chen_sim.py}}).

Bias is the mean of the 1,000 estimates in a cell minus the population value. Each 95\% interval is the 2.5th and 97.5th percentile of 500 bootstrap estimates \citep{efron1993}, with the same resample indices for the three cutoffs, as in her Appendix B. Rejection means the interval excludes zero (Type I error when the population effect is zero, power otherwise); it does not use population values. Coverage is the share of intervals containing the population value, and imbalance splits the misses into the value lying above or below the interval. Bradley's ranges are [.025, .075] for Type I error and [.925, .975] for coverage \citep{bradley1978}; her bias bands are [-.05, .05] and [-.10, .10] \citep{hoogland1998}.

\subsection{Two engines}\label{sec:engines}

The question here is which software produced which number; the answer is that her own printed R code produced the datasets and point estimates it generates, and our Python re-implementation produced every interval-based result.

\textbf{Exact R replay of her seeded code.} Her Appendix A seeds the design table (\texttt{set.seed(18374788)}) and the data (\texttt{set.seed(95409)}). \href{https://github.com/alectionary/vibe-math/blob/master/formalization/chen-rdd-mediation/r-reproduction/chen_exact.R}{\nolinkurl{r-reproduction/chen_exact.R}} runs her own \texttt{gendata()} (defined at line 49) in her order and fits her own Appendix B \texttt{lm()} models, so its datasets and point estimates are exactly those her printed, seeded code generates. They are not necessarily the data behind her published results: on the evidence of Section~\ref{dis:h-absent}, inferred from her figures, those lacked the interaction wherever $h > 0$, which the printed code does not produce. The replay therefore reproduces her printed pipeline, not her published runs. It was run with R 4.1.3 \citep{rcore2022} and packages from the 2021-08-02 snapshot her code pins (dplyr 1.0.7, tidyr 1.1.3, purrr 0.3.4, data.table 1.14.0, stringi 1.7.3, digest 0.6.27); she used R 4.0.3 \citep{rcore2021}, which has the same default random number generator. A literal check runs her own \texttt{\detokenize{group_walk}}/\texttt{\detokenize{pmap_dfr}} pipeline on the first simulation block of her code (1,270 datasets) and compares it with the replay's method: by design the check rebuilds every second dataset (635) and skips the rest by drawing the same number of normal deviates. The 635 rebuilt datasets match her printed pipeline exactly (largest absolute difference 0), and the random-number state after all 1,270 is identical (\href{https://github.com/alectionary/vibe-math/blob/master/formalization/chen-rdd-mediation/r-reproduction/out-full/literal_check.txt}{\nolinkurl{r-reproduction/out-full/literal_check.txt}}). The replay first rebuilt 64,000 datasets (50 replications per condition; \href{https://github.com/alectionary/vibe-math/blob/master/formalization/chen-rdd-mediation/r-reproduction/out/estimates.csv}{\nolinkurl{r-reproduction/out/estimates.csv}}), then all 1,280,000 (1,000 per condition) in 3 hours 25 minutes, writing the mean and standard deviation of every estimate per condition (\href{https://github.com/alectionary/vibe-math/blob/master/formalization/chen-rdd-mediation/r-reproduction/out-full/condition_summary.csv}{\nolinkurl{r-reproduction/out-full/condition_summary.csv}}; run record \href{https://github.com/alectionary/vibe-math/blob/master/formalization/chen-rdd-mediation/r-reproduction/out-full/RUN.md}{\nolinkurl{r-reproduction/out-full/RUN.md}}). It does not reproduce her bootstrap, which sets no seed.

\textbf{Python run for bootstrap-based results.} The full simulation with 500 bootstrap resamples per analysis ran in \href{https://github.com/alectionary/vibe-math/blob/master/formalization/chen-rdd-mediation/simulation/chen_sim.py}{\nolinkurl{simulation/chen_sim.py}}, which computes the same estimators from Gram matrices; bootstrap resamples reuse the data through resample counts. Its estimator was validated by \href{https://github.com/alectionary/vibe-math/blob/master/formalization/chen-rdd-mediation/simulation/validate_estimator.py}{\nolinkurl{simulation/validate_estimator.py}} against scikit-learn 1.9.1 least-squares fits \citep{pedregosa2011} (agreement to $1.1\times 10^{-14}$) and against explicit row resampling. Condition $k$ uses \texttt{\detokenize{SeedSequence(20211201, spawn_key=(k,))}}, our own seed. The run took 13.7 core-hours with Python 3.13.18 and NumPy 2.5.3 \citep{harris2020} on one host and Python 3.10.12 and NumPy 2.2.6 on another; two conditions rerun across hosts agreed to $1.1\times 10^{-16}$ (\href{https://github.com/alectionary/vibe-math/blob/master/formalization/chen-rdd-mediation/notes/run-record.md}{\nolinkurl{notes/run-record.md}}).

\textbf{How the engines were chosen.} The Python run came first. It was chosen because R could not be installed in our computing environment, and that choice was made without first asking the owner of this replication project, although a departure from the source's own software is the owner's decision. When told, the owner agreed to keep Python for the bootstrap-based results and asked for the exact replay of Chen's code in R, which was then run on a separate R-capable machine; the two runs (50 and then 1,000 replications per condition) are documented in the project's public issue tracker as issues \href{https://github.com/alectionary/vibe-math/issues/92}{\#92} and \href{https://github.com/alectionary/vibe-math/issues/118}{\#118}. A full R rerun including the bootstrap (about 4 billion \texttt{lm()} fits) was declined, since an unseeded bootstrap cannot match hers draw for draw in any language. The decision is recorded in \href{https://github.com/alectionary/vibe-math/blob/master/formalization/chen-rdd-mediation/simulation/fallback.json}{\nolinkurl{simulation/fallback.json}}. Consequently bias and point estimates are compared on the datasets her printed code generates, while coverage, imbalance, Type I error and power agree with hers only in distribution.

\subsection{Digitizing her figures}\label{sec:digitize}

Chen prints most of her results only as figures, so we read them. Figures 4.12-4.16 (bias) are vector drawings: every marker was read from the PDF drawing paths, calibrated on her $\pm$.05 and $\pm$.10 reference lines, with reading error at most .0021 (\href{https://github.com/alectionary/vibe-math/blob/master/formalization/chen-rdd-mediation/secondary/digitize_bias_figures.py}{\nolinkurl{secondary/digitize_bias_figures.py}}, output \href{https://github.com/alectionary/vibe-math/blob/master/formalization/chen-rdd-mediation/results/digitized_bias.csv}{\nolinkurl{results/digitized_bias.csv}}). Figures 4.17-4.29 and the odd-numbered coverage Figures 4.31-4.37 are embedded raster images; each panel's axes were calibrated from her reference lines and tick marks, and each series was located by its color at the five sample sizes (\href{https://github.com/alectionary/vibe-math/blob/master/formalization/chen-rdd-mediation/secondary/digitize_raster_figures.py}{\nolinkurl{secondary/digitize_raster_figures.py}}, output \href{https://github.com/alectionary/vibe-math/blob/master/formalization/chen-rdd-mediation/results/digitized_rates.csv}{\nolinkurl{results/digitized_rates.csv}}). Reading errors are up to about .02 on Type I error and power, and on coverage have a median of about .03 and reach about .05 (Figures 4.31 and 4.35). A series hidden under the others at some point is not read there: the digitizer fills that point in from the visible series and flags it as occluded. Such filled-in points are not readings and are not used for statement verdicts, figure-agreement rates or E4. The imbalance bar charts (even-numbered Figures 4.30-4.38) were read from bar heights (\href{https://github.com/alectionary/vibe-math/blob/master/formalization/chen-rdd-mediation/secondary/digitize_imbalance_figures.py}{\nolinkurl{secondary/digitize_imbalance_figures.py}}, output \href{https://github.com/alectionary/vibe-math/blob/master/formalization/chen-rdd-mediation/results/digitized_imbalance.csv}{\nolinkurl{results/digitized_imbalance.csv}}). Coverage is plotted twice, as lines and as bars; the two readings agree within their combined reading error for 898 of 902 directly read points (978 of 982 including merged panels, where series could not be separated; 1,282 of 1,357 if filled-in hidden points are included) (\href{https://github.com/alectionary/vibe-math/blob/master/formalization/chen-rdd-mediation/results/digitize_crosscheck.csv}{\nolinkurl{results/digitize_crosscheck.csv}}). The digitizing scripts use PyMuPDF and pandas; their versions were not recorded with the outputs.

\subsection{Reading her text statements}\label{sec:statements}

Every Type I error and power statement in her Chapter 4 that names cells and a checkable outcome was tested on every cell set it names, with rates averaged over the factors it does not name, as in her figures (\href{https://github.com/alectionary/vibe-math/blob/master/formalization/chen-rdd-mediation/simulation/statement_checks.py}{\nolinkurl{simulation/statement_checks.py}}, output \href{https://github.com/alectionary/vibe-math/blob/master/formalization/chen-rdd-mediation/results/statement_checks.csv}{\nolinkurl{results/statement_checks.csv}}). This gives 18 Type I error and 28 power statements. Her words were given fixed numeric readings (Section~\ref{sec:ass-words}); a second reading moves every threshold .01 in her favor. Each statement is also read against her own digitized figure points and, as a consistency check, against our cells with $h$ replaced by 0.

\subsection{Comparability of the two sides}\label{sec:comparability}

Table~\ref{tab:comparability} records, for her published results and our two engines, the choices that could change a number. By design the sides differ only in random draws for interval-based results and in the step that compares estimates with population values, which she did not print; on the evidence of Section~\ref{dis:h-absent} (an inference from her figures), the data behind her published results also lacked $h$ where $h > 0$.

\begin{table}[ht]
\centering\footnotesize
\caption{Comparability record. ``Hers'' is from her Chapter 3 and Appendices A-B \citep{chen2021}.}\label{tab:comparability}
\begin{tabular}{@{}>{\raggedright\arraybackslash}p{0.17\textwidth}>{\raggedright\arraybackslash}p{0.27\textwidth}>{\raggedright\arraybackslash}p{0.25\textwidth}>{\raggedright\arraybackslash}p{0.23\textwidth}@{}}
\toprule
Item & Hers (published) & Ours, Python run & Her code, R replay\\
\midrule
Data & Appendix A, seeds 18374788 and 95409; published results evidently lack $h$ where $h > 0$ (inferred, Section~\ref{dis:h-absent}) & Same model, our seeds & Datasets her printed code generates; literal check on her first simulation block: 635 of 1,270 datasets rebuilt by design, all identical, random-number state identical after the block; the full replay keeps all 1,280,000\\
Estimator & OLS, $m \sim t + r$, $y \sim t + r + m + tm$ & Same & Her \texttt{lm()} calls\\
$n$ and replications & 5 sizes; 1,000 per condition & Same; 1,000 & Same; 1,000\\
Cutoffs & Code: treated at or above 25th/50th/75th percentile & As her code & As her code\\
Intervals & 500-resample percentile bootstrap, unseeded & 500 resamples, our seeds & Not run\\
Population values & Not printed; PNDE defined as $c' + h$ in text & From the definitions; PNDE $= c'$ & Same as Python run\\
Aggregation & Not printed; collapsed over $n$ and effect sizes & Same collapse (Tables 4.8-4.9) & Condition means\\
\bottomrule
\end{tabular}
\end{table}

\section{Results}\label{sec:results}

\subsection{H1: direct effects are not biased}\label{sec:h1}

Are the natural direct effects biased, as \citet[Table 4.8, printed p.~88; Chapter 5, p.~147]{chen2021} concluded? No. In our run every PNDE and TNDE result lies within [-.05, .05]; the largest $|$bias$|$ over all 19,200 condition-cutoff-effect results, for all five effects, is 0.016, and the Monte Carlo standard error of a cell mean is at most 0.006 (\href{https://github.com/alectionary/vibe-math/blob/master/formalization/chen-rdd-mediation/results/table48_ours.csv}{\nolinkurl{results/table48_ours.csv}}, \href{https://github.com/alectionary/vibe-math/blob/master/formalization/chen-rdd-mediation/results/summary.json}{\nolinkurl{results/summary.json}}). She reported that only 12.50-18.75\% of direct-effect results fall within [-.10, .10] and that the bias is negative. The signed contrast E1 (Table~1) is 85.42 points in her table and 0.00 in ours.

Her printed code agrees with our run. In the exact replay over all 1,000 replications, the largest condition-mean bias is 0.017 and none exceeds .10; 5.3\% of the 19,200 cell biases lie beyond 1.96 Monte Carlo standard errors and 0.14\% beyond 3.29, as expected with no bias; her estimates and ours differ beyond 1.96 combined standard errors in 5.3\% of cells; the median ratio of their standard deviations is 1.000; and the slope of estimated on true $h$ is 0.9998-1.0005 by cutoff (\href{https://github.com/alectionary/vibe-math/blob/master/formalization/chen-rdd-mediation/results/r_full_comparison.json}{\nolinkurl{results/r_full_comparison.json}}, per cell in \href{https://github.com/alectionary/vibe-math/blob/master/formalization/chen-rdd-mediation/results/r_full_comparison.csv}{\nolinkurl{results/r_full_comparison.csv}}). Her digitized bias (Figures 4.12-4.16) is instead an almost exact function of the design with no dependence on $n$: for PNDE, $-1.466$, $-0.965$ and $-0.446$ times \cp{} at cutoffs 1-3, plus about $-0.928\,h$; for MINT, $-1.00\,ah$ (\href{https://github.com/alectionary/vibe-math/blob/master/formalization/chen-rdd-mediation/results/bias_structure.csv}{\nolinkurl{results/bias_structure.csv}}). Sections~\ref{dis:h-absent}, \ref{dis:population-values} and \ref{dis:pnde-definition} give the explanations and their tests. H1 is not replicated.

\subsection{H2: percentile intervals have close to nominal coverage}\label{sec:h2}

Do 95\% percentile bootstrap intervals cover the population value too rarely, with the value systematically above the interval \citep[Table 4.9, printed p.~124; Chapter 5, p.~150]{chen2021}? Not in our run. Mean coverage is 94.6-95.6\% in every effect-cutoff row, against her 24.10-83.99\%; the lowest cell coverage is 0.916, and 92.9\% of the 19,200 results lie in Bradley's [.925, .975] (\href{https://github.com/alectionary/vibe-math/blob/master/formalization/chen-rdd-mediation/results/table49_ours.csv}{\nolinkurl{results/table49_ours.csv}}). ``Close to nominal'' describes the row means: the cells outside Bradley's range are concentrated in the product effects, with 80.1\% of MINT cells, 89.4\% of PNIE and 95.4\% of TNIE cells inside it, against 99.8\% for PNDE and TNDE (our calculation from \href{https://github.com/alectionary/vibe-math/blob/master/formalization/chen-rdd-mediation/results/cells.csv}{\nolinkurl{results/cells.csv}}). Misses fall on both sides at similar rates: the value lies below the interval in 1.8-2.7\% of replications and above it in 2.5-2.9\%. With 1,000 replications a coverage near .95 has a binomial standard error of about .007. Contrast E2, 95 minus mean coverage, is 43.62 points in her table and 0.04 in ours; contrast E3, the mean share above minus the share below, is 47.07 and 0.33.

Our run shows a small positive imbalance (0.33 points, against her 47.07), about two orders of magnitude smaller than hers; we did not investigate its cause. Her intervals sit where ours do: her PNIE Type I error and power agree with ours at 100.0\% and 97.2\% of directly read points (\href{https://github.com/alectionary/vibe-math/blob/master/formalization/chen-rdd-mediation/results/rates_agreement.csv}{\nolinkurl{results/rates_agreement.csv}}). Her digitized coverage of nonzero PNIE falls with $n$ to a plateau of 0.25, 0.487 and 0.736 at cutoffs 1-3, the control share under her code, with the value above the interval in about the treated share (Section~\ref{dis:population-values}). H2, and her conclusion that the percentile interval does not give valid inference for mediation in RDD, are not replicated.

\subsection{H3: MINT power rises with sample size}\label{sec:h3}

Is power for the mediated interaction near zero at every sample size \citep[Figure 4.28, printed p.~122; text p.~121; Chapter 5, p.~148]{chen2021}? No. In our run MINT power rises with $n$ and with $ah$: 0.836 at $n = 500$ for $ah = .2006$ (the largest product in the design), 0.719 at $n = 500$ for $ah = .1326$, and 0.9735 or more for $ah \ge .13$ from $n = 1{,}000$ (\href{https://github.com/alectionary/vibe-math/blob/master/formalization/chen-rdd-mediation/results/summary.json}{\nolinkurl{results/summary.json}}, \href{https://github.com/alectionary/vibe-math/blob/master/formalization/chen-rdd-mediation/results/power.csv}{\nolinkurl{results/power.csv}}). Her Figure 4.28 reads 0.068 and 0.057 in the contrast cell at cutoffs 2 and 3 (unrounded .0683 and .0571, whose mean .0627 gives the 0.063 and 0.0127 used below), and no point above about .07, which is the shape of a Type I error curve; her cutoff-1 series is hidden under the others there and could not be read, so E4 uses cutoffs 2-3 on both sides. Contrast E4 (Table~1) therefore measures each rate against its null: the rejection rate minus the nominal .05 is 0.0127 in her figure (raw 0.063) and 0.802 in our run (raw 0.852, Monte Carlo standard error about .002). Her excess is about 1.4 reading errors: small and positive, and far below ours, so the magnitude is larger in our run and the direction is the same. The published value is read from an image; the tolerance, 0.0091, is the mean reading error of her two read points, rounded down. Over all three cutoffs our rate in this cell is 0.836. This contrast cell, the largest MINT in the design at the smallest $n$, was chosen after the full results were available; \href{https://github.com/alectionary/vibe-math/blob/master/formalization/chen-rdd-mediation/results/summary.json}{\nolinkurl{results/summary.json}} gives MINT power for every $n$ and $ah$.

Rejection rates do not use population values, so this disagreement cannot come from the comparison with population values. Her MINT power matches our rates computed with $h$ set to 0 at 96.8\% of directly read points, against 4.3\% with $h$ as designed, and her MINT bias is $-1.00\,ah$; her printed code, replayed, recovers $h$ (mean MINT estimate 0.2009 against a population value of 0.2006 in this cell). Her own summaries are also inconsistent on this point: her abstract and Chapter 5 (printed p.~143) say that Type I error and power for all five effects were within or reached the robustness criteria, which her MINT power statement contradicts; we test the specific Chapter 4-5 statement because it names the effect and the outcome. H3 is not replicated (Section~\ref{dis:h-absent}).

\subsection{H4: the 5,000-case recommendation holds for most cells}\label{sec:h4}

Does $n = 5{,}000$ give accurate Type I error and power above .80 regardless of cutoff \citep[Chapter 5, printed p.~150]{chen2021}? Mostly. Chen states this as tentative; we judge it literally over the whole grid. The share of zero-effect results with Type I error in Bradley's range rises from 70.0\% at $n = 500$ to 74.0\%, 80.1\%, 86.3\% and 89.6\% at $n = 1{,}000$ to 10,000, and the share of nonzero results with power of at least .80 from 35.9\% to 53.3\%, 61.5\%, 89.2\% and 91.4\%; the largest step is from 2,500 to 5,000 (\href{https://github.com/alectionary/vibe-math/blob/master/formalization/chen-rdd-mediation/results/statement_checks.csv}{\nolinkurl{results/statement_checks.csv}}, over 1,176 zero-effect and 2,664 nonzero results per $n$). At $n = 5{,}000$ the shares differ little by cutoff (85.7-87.0\% and 88.3-89.9\%). Not every result passes: product effects whose two paths are both zero stay conservative (about .002) at every $n$, in her Figure 4.17 as in our run, and small effects such as $ah = .007$ stay below .80 power (0.328 at $n = 5{,}000$ and 0.669 at 10,000). H4 is therefore partly replicated (Section~\ref{dis:n5000-overgeneralized}).

\subsection{S1: indirect and MINT bias}\label{sec:s1}

\citet[Table 4.8]{chen2021} reported that indirect and MINT bias falls within [-.10, .10] in 51.88-87.50\% of conditions, more often as the control share grows. Our indirect and MINT results all lie within [-.05, .05], in our run and in the replay of her code; the share outside [-.10, .10], averaged over cutoffs, is 27.08\% (PNIE), 39.06\% (TNIE) and 12.50\% (MINT) in her table and 0.0\% for each in ours. Her digitized PNIE bias is $-k\,ab$ with $k = 1.469$, 0.959 and 0.441 at cutoffs 1-3 and no dependence on $n$, and her MINT bias is $-ah$ (\href{https://github.com/alectionary/vibe-math/blob/master/formalization/chen-rdd-mediation/results/bias_structure.csv}{\nolinkurl{results/bias_structure.csv}}). At the median cutoff, her PNIE and MINT rows and her TNIE [-.05, .05] entry also equal what an estimate of zero would give over the grid (\href{https://github.com/alectionary/vibe-math/blob/master/formalization/chen-rdd-mediation/secondary/zero_estimate_check.csv}{\nolinkurl{secondary/zero_estimate_check.csv}}); that pattern does not fit every printed entry (PNDE, for example, needs \cp{} + $h$), so the bias-structure fit is the fuller description. S1 is not replicated.

\subsection{S2: Type I error statements}\label{sec:s2}

Of her 18 testable Type I error statements (Figures 4.17-4.21; printed pp. 99-108), 13 agree with our run and 14 with .01 threshold slack. The five that disagree as designed are statements 2, 9, 10, 11 and 18 (the \texttt{number} column of \href{https://github.com/alectionary/vibe-math/blob/master/formalization/chen-rdd-mediation/results/statement_checks.csv}{\nolinkurl{results/statement_checks.csv}}). Statement 2 agrees with slack (Section~\ref{dis:threshold-noise}). Statement 18 concerns MINT with $a = 0$: our rate is below Bradley's range at $h = .05$ through $n = 5{,}000$ (at most .019) and at $h = .15$ at $n = 500$ (.014-.020), as she reports, but it rises with $n$ at $h = .15$ (.026-.033 at $n = 1{,}000$, about .05 from $n = 2{,}500$) and is about .05 at $h = .34$, whereas she reports near-zero rates with no variation across $h$. Statements 9, 10 and 11 concern TNIE with $a = 0$ at small $n$, where our rates and her description fall on different sides of the bounds. Read on our cells with $h = 0$, statements 9 and 18 hold; statements 10 and 11 disagree under every reading of our run, and statement 10 sits within one standard error of the .025 threshold (Section~\ref{dis:threshold-noise}). The numeric readings of her words are ours, and points hidden by another series are excluded. S2 is partly replicated.

\subsection{S3: power statements}\label{sec:s3}

Of her 28 testable power statements (Figures 4.22-4.28; printed pp. 110-121), 16 agree with our run and 17 with slack. All 6 PNIE statements and all 5 PNDE statements agree; the disagreements concern TNIE, TNDE and MINT, mostly in cells in which $h$ is part of the effect. Over all 46 statements, 29 agree as designed and 31 with slack. Read on our cells with $h$ replaced by 0, 39 of 46 hold as designed (strict reading); 10 of the 15 statements that disagree under both readings become agreements, and none that agrees becomes a disagreement. The 7 that disagree with $h = 0$ are statements 2 and 43, which agree with the .01 slack (Section~\ref{dis:threshold-noise}), and statements 10, 11, 30, 41 and 44, which disagree under every reading (the \texttt{number} column of \href{https://github.com/alectionary/vibe-math/blob/master/formalization/chen-rdd-mediation/results/statement_checks.csv}{\nolinkurl{results/statement_checks.csv}}). Her digitized Figures 4.17-4.28 match our rates read with $h = 0$ at 79.1-100.0\% of directly read points per figure, against 4.3-100.0\% with $h$ as designed (\texttt{\detokenize{*_read}} columns of \href{https://github.com/alectionary/vibe-math/blob/master/formalization/chen-rdd-mediation/results/rates_agreement.csv}{\nolinkurl{results/rates_agreement.csv}}). One figure goes the other way: Figure 4.25 (PNDE power) agrees at 85.9\% with $h$ as designed and 79.1\% with $h = 0$. The remaining disagreements are threshold cases, her text against her own figures, and one TNDE cell set explained by her particular random draws (Sections~\ref{dis:threshold-noise}, \ref{dis:prose-versus-own-figure} and \ref{dis:tnde-power-gap}). S3 is partly replicated.

\subsection{RQ2: the Head Start analysis (S4, S5, S6)}\label{sec:rq2}

Can her comparison of randomized and synthetic-RDD mediation estimates in the Head Start Impact Study be rerun? No: it uses the restricted-use HSIS file (3,721 children, 2,697 with a wave-1 spelling score), which requires a signed ICPSR agreement this project does not hold (Section~\ref{sec:availability}). Her conclusion that RDD estimates are typically smaller than the trial estimates, with standardized bias outside an acceptable range (S4; Tables 4.6-4.7, printed pp. 145-146), and her assumption checks showing slight violations of independence and functional form (S6; Tables 4.3-4.5) are therefore blocked: not established, and not a failure to replicate.

What can be checked from her printed numbers holds (S5; \href{https://github.com/alectionary/vibe-math/blob/master/formalization/chen-rdd-mediation/secondary/table_arithmetic.csv}{\nolinkurl{secondary/table_arithmetic.csv}}). All 24 printed RCT-minus-RDD differences in Tables 4.6-4.7 recompute. All 24 standardized biases match only when the divisor is the standard error of the RCT estimate, not the standard deviation of the RCT outcome stated in her equation 3.12 (Section~\ref{dis:headstart-sb-divisor}); with that divisor the printed sign is that of RDD minus RCT. Her printed effects obey their definitions: 12 of 12 identities hold (PNIE $= ab$, TNIE $= ab + ah$ and TNDE $-$ PNDE $= ah$ in the trial and at three cutoffs; \href{https://github.com/alectionary/vibe-math/blob/master/formalization/chen-rdd-mediation/results/empirical_consistency.json}{\nolinkurl{results/empirical_consistency.json}}), so on real data her code estimated $h$ and used it. The PNDE identity implies a mediator intercept of 19.67 (trial) and 22.11 (cutoff 3), the scale of the mediator mean of about 21 minutes in her Table 4.2. In Table 4.2 treatment is below the cutoff, with treated shares .356, .717 and .856 at cutoffs 1-3, as her text labels them and the reverse of her simulation code. In Table 4.3, 8 of 9 $t$ values equal the coefficient over its standard error (\href{https://github.com/alectionary/vibe-math/blob/master/formalization/chen-rdd-mediation/results/table_identities.json}{\nolinkurl{results/table_identities.json}}; Section~\ref{dis:table-4-3-t-value}). S5 is replicated at the printed precision.

Her abstract and Chapter 5 (printed p.~143) also conclude that 95\% intervals for the RDD estimates reach opposite significance conclusions to the trial estimates: if the trial estimate was significant, the RDD estimate was not (S9). Read from her own printed Table 4.7 intervals (\href{https://github.com/alectionary/vibe-math/blob/master/formalization/chen-rdd-mediation/source-values/table_4_7.csv}{\nolinkurl{source-values/table_4_7.csv}}; \href{https://github.com/alectionary/vibe-math/blob/master/formalization/chen-rdd-mediation/results/significance_conclusions.json}{\nolinkurl{results/significance_conclusions.json}}, code \href{https://github.com/alectionary/vibe-math/blob/master/formalization/chen-rdd-mediation/secondary/significance_conclusions.py}{\nolinkurl{secondary/significance_conclusions.py}}), 10 of the 15 trial-versus-RDD comparisons are opposite and 5 agree (PNDE and TNDE at cutoff 3, MINT at all three cutoffs). Where the trial estimate is significant (PNIE and TNIE), the RDD estimate is not in all 6 cells, as she says; but in 4 cells (PNDE and TNDE at cutoffs 1-2) the RDD interval excludes zero where the trial interval does not, so her printed p.~81 sentence that PNDE and TNDE ``were not significant'' holds for the trial only. Her own Chapter 4 reading of Table 4.7 (printed p.~85) agrees with ours cell by cell: it reports the RDD intervals for PNIE and TNIE as Type II errors and those for PNDE and TNDE at cutoffs 1-2 as Type I errors; only her abstract and the p.~143 summary overgeneralize. S9 is partly replicated, from printed numbers only (Section~\ref{dis:significance-conclusions}).

\subsection{RQ1 and the announced power analysis (S7, S8)}\label{sec:s7s8}

Her RQ1 answer is an integrated table of mediation and RDD assumptions (S7; Table 4.1, printed p.~60). It is a literature synthesis with no computation to rerun, so it is not established here and was not evaluated as scholarship; we note only that the simulation model meets the no-confounding and functional-form assumptions by construction. Her Chapter 3 announces an additional power analysis comparing expected and simulated power by the standard deviation of $Y$ (S8; printed p.~58), but it was never reported: no method or result appears in Chapters 4-5 or the appendices (searched for ``expected power'', ``standard deviation'' and ``power analysis''), so there is nothing to assess: S8 is not established, and nothing blocks it.

\section{Comparison of published and replicated values}\label{sec:comparison}

Table~\ref{tab:comparison} lists each published quantity with our value. Numerical is \emph{match} when our value rounds to the printed value, \emph{close} when it misses by at most one unit in the last printed place, and \emph{mismatch} otherwise; scope is \emph{exact} for the same data, estimator and procedure, \emph{limited} for a reconstruction (here, the same model with our own random draws, or printed numbers only), and \emph{unavailable} when the input was not obtained. \emph{Withheld} means no numerical comparison is possible because the input was not obtained; \emph{not applicable} means the claim has no number to compare (a qualitative, threshold or conceptual claim); the last column compares signs, and \emph{not established} means no signed comparison exists.

{\footnotesize
\setlength{\tabcolsep}{3pt}
\begin{longtable}{@{}>{\raggedright\arraybackslash}p{0.16\textwidth}>{\raggedright\arraybackslash}p{0.15\textwidth}>{\raggedright\arraybackslash}p{0.16\textwidth}>{\raggedright\arraybackslash}p{0.11\textwidth}>{\raggedright\arraybackslash}p{0.10\textwidth}>{\raggedright\arraybackslash}p{0.10\textwidth}>{\raggedright\arraybackslash}p{0.12\textwidth}@{}}
\caption{Published values \citep{chen2021} against ours. Reasons for each row are in the paragraph that follows.}\label{tab:comparison}\\
\toprule
Claim & Published & Ours & Precision & Numerical & Scope & Same direction?\\
\midrule\endfirsthead
\toprule
Claim & Published & Ours & Precision & Numerical & Scope & Same direction?\\
\midrule\endhead
H1 direct effects, \% outside [-.10, .10] & 85.42 & 0.00 & 2 decimals & mismatch & limited (same data at $h = 0$) & zero vs nonzero\\
H2 coverage shortfall (points) & 43.62 & 0.04 & 2 decimals & mismatch & limited & yes\\
H2 above minus below (points) & 47.07 & 0.33 & 2 decimals & mismatch & limited & yes\\
H3 MINT power minus .05, $n = 500$, $ah = .20$, cutoffs 2-3 & 0.0127 (raw 0.063, figure) & 0.802 (raw 0.852) & reading error .0091 & mismatch & limited & yes\\
H4 $n = 5{,}000$ recommendation & accurate Type I and power $>$ .80 & 86.3\% and 89.2\% of results & qualitative & not applicable & limited & not established\\
S1 PNIE \% outside [-.10, .10] & 27.08 & 0.00 & 2 decimals & mismatch & limited (same data at $h = 0$) & zero vs nonzero\\
S1 TNIE \% outside & 39.06 & 0.00 & 2 decimals & mismatch & limited (same data at $h = 0$) & zero vs nonzero\\
S1 MINT \% outside & 12.50 & 0.00 & 2 decimals & mismatch & limited (same data at $h = 0$) & zero vs nonzero\\
S2 Type I statements & 18 statements & 13 agree (14 with slack) & thresholds & not applicable & limited & not applicable\\
S3 power statements & 28 statements & 16 agree (17 with slack) & thresholds & not applicable & limited & not applicable\\
S4 HSIS RDD versus RCT & RDD smaller & not rerun & --- & withheld & unavailable & not established\\
S5 HSIS arithmetic & 24 differences and 24 standardized biases & all 48 reproduced; 12 of 12 identities hold & 3 decimals & match & limited & not applicable\\
S6 HSIS assumption checks & slight violations & not rerun & --- & withheld & unavailable & not established\\
S7 assumption table & Table 4.1 & conceptual & --- & not applicable & not applicable & not applicable\\
S8 power by SD of $Y$ & announced, never reported & not established & --- & not applicable & not applicable & not established\\
S9 HSIS significance conclusions & opposite in every comparison & 10 of 15 opposite & printed interval bounds & not applicable & limited & not applicable\\
\bottomrule
\end{longtable}}

Reasons by row. H1 and S1: the replay reproduces the datasets her printed code generates, and it gives 0\% outside the band. Her published runs lacked the interaction wherever $h > 0$ (Section~\ref{dis:h-absent}), so the replay is not her published runs and the scope is limited; it is exact only where $h = 0$, where both candidate mechanisms for the missing interaction leave the data unchanged. The mismatch stands there too, because her bias pattern ($-k\,ab$ for PNIE, $-k\,c'$ for PNDE) does not depend on $h$. H2: our coverage rows are within a few tenths of a point of 95\% while hers are 11-71 points below it; scope is limited because intervals use our draws. H3: against the nominal .05, her rate is about 1.4 reading errors above it (0.0127 against a tolerance of 0.0091) while ours is 0.802 above it; her value is the mean of two figure readings and ours a 32-cell average. H4: a recommendation, judged by shares over the grid. S2-S3: threshold statements, judged per cell set. S4 and S6: inputs unavailable. S5: all printed differences and standardized biases recompute (with the estimate's standard error as divisor). S7: conceptual. S8: announced but never reported, so there is nothing to compare. S9: significance read from her printed intervals; the fits cannot be rerun.

\textbf{Small differences that are not replication failures.} (1) One Table 4.3 $t$ value disagrees with its coefficient and standard error, a one-digit typo that changes no conclusion (Section~\ref{dis:table-4-3-t-value}). (2) The standardized-bias divisor in her tables is not the one her equation 3.12 states (Section~\ref{dis:headstart-sb-divisor}). (3) Four Table 4.9 rows sum to 99.99 or 100.01 rather than 100, which is rounding (\href{https://github.com/alectionary/vibe-math/blob/master/formalization/chen-rdd-mediation/results/table_identities.json}{\nolinkurl{results/table_identities.json}}). (4) Her cutoff labels in the text are the reverse of her simulation code (Section~\ref{sec:amb-cutoff}).

\section{Disagreements}\label{sec:disagreements}

Each disagreement with \citet{chen2021} was investigated with competing explanations, at least one of them an error of ours, which was tested first. The full record, with the files behind each test, is \href{https://github.com/alectionary/vibe-math/blob/master/formalization/chen-rdd-mediation/DISAGREEMENTS.yaml}{\nolinkurl{DISAGREEMENTS.yaml}}; most test outputs are in \href{https://github.com/alectionary/vibe-math/blob/master/formalization/chen-rdd-mediation/results/disagreement_tests.json}{\nolinkurl{results/disagreement_tests.json}} (code \href{https://github.com/alectionary/vibe-math/blob/master/formalization/chen-rdd-mediation/secondary/disagreement_tests.py}{\nolinkurl{secondary/disagreement_tests.py}}). Of the ten, eight are explained and two (Sections~\ref{dis:population-values} and \ref{dis:pnde-definition}) are partly explained; none is unexplained.

\subsection{Effects involving the interaction behave as if it were absent}\label{dis:h-absent}

\textbf{Her result.} Effects that involve the treatment-by-mediator interaction $h$ behave as if $h$ were absent: MINT bias is $-ah$ at every cutoff (Figure 4.16), MINT power stays near .05 (Figure 4.28), and Type I error with $a = b = 0$ and $h > 0$ is near .001 (Figures 4.18 and 4.21); TNIE and TNDE carry the same pattern. \textbf{Ours.} $h$ is recovered: mean MINT bias by $ah$ is within .001 and every cell within .007, MINT power reaches .84 at $n = 500$ and at least .97 from $n = 1{,}000$ for $ah \ge .13$, and those Type I rates are near .05.

\begin{description}[leftmargin=1em,style=nextline]
\item[Our implementation loses $h$ (ours).] Prediction: her printed code, run on the datasets its seeds generate, would behave like her results. Test: the exact replay of her Appendix A \texttt{gendata()} and Appendix B \texttt{lm()} calls, first on 64,000 datasets, then on all 1,280,000. Result: her code recovers $h$ (slope of estimated on true $h$ 0.9998, 1.0005 and 0.9999 at the three cutoffs); the largest condition-mean bias for any effect is 0.017; 5.3\% of cell biases lie beyond 1.96 Monte Carlo standard errors and 0.14\% beyond 3.29; her estimates and ours differ beyond 1.96 combined standard errors in 5.3\% of cells, with a median standard-deviation ratio of 1.000. Verdict: ruled out.
\item[We misread her figures (ours).] Prediction: the pattern would not survive an exact reading, and her two coverage plots would disagree. Test: vector reading of Figures 4.12-4.16 calibrated on her reference lines; the double reading of coverage. Result: the MINT bias fit is $-1.00\,ah$ with residual at most .001 against a reading error of at most .0021; the coverage readings agree for 898 of 902 directly read points (978 of 982 with merged panels; 1,282 of 1,357 with filled-in hidden points). Verdict: ruled out.
\item[Her printed code drops $h$ (source).] Prediction: the replay would show $\hat h$ near 0. Test and result: the same replay recovers $h$. Verdict: ruled out.
\item[Her fitted $y$ model omitted $tm$ (source).] Prediction: $\hat b$ would absorb the interaction, so her PNIE bias would carry a positive $ah$ term growing with the treated share. Test: least-squares fit of her digitized PNIE bias on design terms, and her real-data estimates. Result: the $ah$ and $h$ coefficients of her PNIE bias are 0.000 at every cutoff (residual at most .005), and on the Head Start data her code estimated $h$ (.003, .001, .001, $-$.028) and used it, since TNIE $= ab + ah$ holds in all four designs. Verdict: ruled out.
\item[The data behind her results had no interaction (source).] Prediction: her rejection rates would match ours read with $h = 0$, and her MINT estimates would average 0. Test: every digitized Type I and power point and every statement compared with our cells with $h$ replaced by 0. Result: 79.1-100.0\% of directly read points agree in every Type I and power figure (MINT power from 4.3\% to 96.8\%); filled-in hidden points are not counted; 10 of the 15 disagreeing statements agree; MINT bias is $-1.00\,ah$ at all three cutoffs. Verdict: supported.
\item[Which mechanism removed the interaction (source).] Two candidates: $y$ generated before the $tm$ column was filled (her \texttt{gendata()} pre-fills the matrix with 1s, so $h$ would multiply 1), or the $h$ value not reaching \texttt{gendata()}. Prediction: they differ only in the $y$ intercept. Test: both mechanisms on one simulated dataset with the same draws. Result: only the intercept differs (by $h$); the $t$, $r$, $m$ and $tm$ coefficients are identical. Her $y$ intercepts are not reported, so her published results cannot separate them. Verdict: inconclusive.
\end{description}
\textbf{Conclusion.} Explained; the error is on the source side, with high confidence: the evidence indicates that the datasets behind her published results contained no treatment-by-mediator interaction (an inference from her outputs; her code and data for those runs are not available), while her printed code recovers it. Which of the two mechanisms did this is not determined.

\subsection{Coverage and bias track the treated share}\label{dis:population-values}

\textbf{Her result.} Coverage of nonzero effects falls with $n$ to a plateau at the control share (.25, .49, .74 at cutoffs 1-3 from $n = 2{,}500$; Figures 4.29-4.38, Table 4.9), with the true value mostly above the interval; the bias of the path products is $-k$ times the product, with $k = 1.47$, .96 and .44 by cutoff and no dependence on $n$ (Figures 4.12-4.15, Table 4.8). \textbf{Ours.} Mean coverage 0.04 points short of 95\% (rows 94.59-95.63\% by effect and cutoff; lowest single cell .916, and 92.9\% of cells inside Bradley's [.925, .975]); every bias within .016 of zero.

\begin{description}[leftmargin=1em,style=nextline]
\item[Our bias or coverage computation is wrong (ours).] Prediction: her printed code would be biased like her results, or her rejection decisions would differ from ours. Test: the exact replay for bias; her digitized PNIE Type I error and power against ours. Result: her code is unbiased, and her PNIE Type I error and power agree with ours at 100\% and 97.2\% of directly read points, so her intervals' locations and widths match ours; only the comparison with population values differs. Verdict: ruled out.
\item[The plateau is a misreading (ours).] Prediction: the line and bar plots of coverage would disagree. Test and result: 898 of 902 directly read points agree within their combined reading error (978 of 982 with merged panels; 1,282 of 1,357 with filled-in hidden points). Verdict: ruled out.
\item[One wrong population value per condition (source).] Prediction: intervals shrink around the estimate as $n$ grows, so coverage would fall toward 0. Test: coverage of nonzero PNIE by cutoff and $n$ in her Figure 4.29. Result: coverage stops at .25, .487 and .736, not 0. Verdict: ruled out.
\item[Population values right for control cases, wrong for treated cases, for coverage (source).] Population values were computed per case, right for control cases and wrong for treated ones, and coverage was averaged over cases. Prediction: coverage falls with $n$ to a plateau near the control share; the share above approaches the treated share and the share below is small. Test: her digitized coverage and imbalance by cutoff and $n$. Result: coverage plateaus at the control share and ``above'' at the treated share at every cutoff, with ``below'' .008-.017. Verdict: supported.
\item[The same rule, with one error scale, for bias (source).] The same per-case rule, with one scale of error for every cutoff, also produced her bias: bias $= -\lambda \times$ treated share $\times$ path product. Prediction: the bias multipliers divided by the treated share give one $\lambda$. Test: the multipliers of her digitized bias. Result: $\lambda = 1.96$, 1.92 and 1.76 at cutoffs 1-3, close for cutoffs 1-2 but not constant, so one scale does not reproduce cutoff 3. Verdict: ruled out; the rule that produced her bias is not determined.
\item[Her treatment labels were reversed in that step (source).] Her text calls cutoff index 1 25\% treated (Chapter 3, Table 3.1, Table 4.8 note), but her code makes it 75\% treated, so a step written for the text's labels would swap treated and control cases. Prediction: the plateau would equal the treated share as her text labels it. Result: the plateau equals both that share and the control share under her code, so it cannot separate this theory from the previous one. The label conflict itself is confirmed: in her Head Start analysis (Table 4.2) treatment is below the cutoff, with treated shares .36, .72 and .86, as her text labels them. Verdict: inconclusive.
\item[Monte Carlo noise (neither).] Prediction: the pattern would vary with $n$, $a$ and $b$ and would not track the treated share. Result: the bias fits leave residuals of at most .024 over 240-960 points per figure, and the plateaus are identical across $a$ and $b$. Verdict: ruled out.
\end{description}
\textbf{Conclusion.} Partly explained; the error is on the source side, with medium confidence. The comparison of estimates with population values, a step she did not print, is where her results depart from her code, and a per-case rule that is right for control cases and wrong for treated ones reproduces her coverage and imbalance at every cutoff. The rule behind her bias is not determined: a single error scale fits cutoffs 1-2 but not cutoff 3.

\subsection{Her definition of the population PNDE}\label{dis:pnde-definition}

\textbf{Her statement.} Her text defines the population PNDE as \cp{} + $h$ (printed p.~104). \textbf{Ours.} Under her data-generating model the population PNDE is $c' + h\,i_{0M}$ with $i_{0M} = 0$, so it is \cp{}.
\begin{description}[leftmargin=1em,style=nextline]
\item[The mediator intercept is not 0 in her design (ours).] Prediction: the mediator-model intercept fitted on her data would be away from 0. Test: mean $\hat\imath_{0M}$ in the exact replay. Result: 0.0001, 0.0000 and 0.0000 at the three cutoffs. Verdict: ruled out.
\item[She used \cp{} + $h$ as the population value (source).] Prediction: with $\hat h$ near 0, her PNDE bias would carry a coefficient of $-1$ on $h$. Test: fit of her digitized PNDE and TNDE bias. Result: the $h$ coefficient is about $-0.93$ (PNDE) and $-1.0$ (TNDE) at every cutoff. The PNDE fit also has an $ah$ coefficient of about $-0.26$ at every cutoff, which this hypothesis does not predict and which is not explained; the $-0.93$ against $-1$ may share that cause. Verdict: supported.
\item[Her estimation code computes PNDE as \cp{} + $h$ in general (source).] Prediction: on her real data, PNDE $-$ \cp{} would equal $h$. Test: the PNDE identity on her printed Head Start estimates. Result: the implied intercept is 19.7 (trial) and 22.1 (cutoff 3), the scale of the mediator mean; at cutoffs 1-2 $h$ rounds to .001, too small to identify it. Her estimation code used the fitted intercept, so the \cp{} + $h$ form belongs to her simulation's population values only. Verdict: ruled out.
\item[Noise (neither).] Prediction: the coefficient would vary across cutoffs with large residuals. Result: it is stable across cutoffs, with residuals at most .010. Verdict: ruled out.
\end{description}
\textbf{Conclusion.} Partly explained; source-side error, medium confidence. Her population PNDE of \cp{} + $h$ accounts for the $h$ term of her bias, but the $ah$ term of about $-0.26$ is not accounted for.

\subsection{The standardized-bias divisor in the Head Start tables}\label{dis:headstart-sb-divisor}

\textbf{Her statement.} Equation 3.12 divides the RCT-minus-RDD difference by the standard deviation of the RCT outcome. \textbf{Ours.} All 24 printed standardized biases equal the difference divided by the standard error of the RCT estimate, with the sign of RDD minus RCT.
\begin{description}[leftmargin=1em,style=nextline]
\item[We transcribed or recomputed her tables wrongly (ours).] Prediction: the printed differences would not recompute either. Test: recompute all 24 differences and standardized biases. Result: all 24 differences match. Verdict: ruled out.
\item[Her computation used the estimate's standard error (source).] Prediction: every printed value would match with that divisor. Result: all 24 match; the implied divisor equals the printed standard error to the printed precision (for example 0.7327-0.7331 against 0.733). Verdict: supported.
\item[Rounding (neither).] Prediction: implied divisors would scatter around the outcome standard deviation. Result: they equal the estimates' standard errors exactly. Verdict: ruled out.
\end{description}
\textbf{Conclusion.} Explained; source-side, high confidence. Which divisor she intended is unknown.

\subsection{Statements on a threshold's knife edge}\label{dis:threshold-noise}

\textbf{Her statements.} Statements 2, 10 and 43 (and some cells of 11 and 41) place a rate on one side of a threshold (.025 or .70, or a .05 spread across cutoffs). \textbf{Ours.} Our rate falls just on the other side.
\begin{description}[leftmargin=1em,style=nextline]
\item[Our rates differ from hers in these cells (ours).] Prediction: our values would sit several standard errors from her digitized values. Test: distance from her value and from the threshold in Monte Carlo standard errors. Result: statement 10 (TNIE at cutoff 3, $a = 0$, $b = .14$, $h = 0$, $n = 1{,}000$): ours .0275, hers .0241, 0.7 standard errors apart and 1.0 from .025; for statement 2 her own figure reads .026, .024 and .024 at $n = 2{,}500$, on both sides of .025. Verdict: ruled out.
\item[Knife edge (neither).] Prediction: moving each threshold .01 in her favor would turn these into agreements. Result: statements 2 and 43 agree with slack; statement 10's $h = 0$ cell is 1 standard error from .025, and its $h = .05$ cell is part of Section~\ref{dis:h-absent}. Verdict: supported.
\end{description}
\textbf{Conclusion.} Explained; neither side's error, medium confidence.

\subsection{Her text against her own figures}\label{dis:prose-versus-own-figure}

\textbf{Her statements.} Statement 30 (TNIE at cutoffs 2 and 3 ``similar to'' cutoff 1), statement 44 (the same for TNDE), and statement 11 (TNIE Type I error ``within'' the range in all other $a = 0$ cases). \textbf{Ours.} Read strictly (no spread across cutoffs above .05; within .025-.075), our run disagrees.
\begin{description}[leftmargin=1em,style=nextline]
\item[Our reading is too strict (ours).] ``Similar patterns'' may mean the same shape over $n$. Prediction: under that reading her figures and our run both agree. Test: whether power rises with $n$ at every cutoff and cell set, in her Figures 4.24 and 4.27 and in our cells. Result: true for all 36 (statement 30) and 48 (statement 44) cell sets, in her figures and ours. Verdict: supported.
\item[Her text contradicts her own figures (source).] Prediction: her digitized figures would fail the strict reading too. Test: every statement read against her own directly read points. Result: her Figure 4.24 spreads by up to .33 across cutoffs at $n = 500$ (ours .34 with $h = 0$); her Figure 4.27 spreads by .051 (ours .046 with $h = 0$); her Figure 4.18 shows $a = b = 0$, $h > 0$ at about .001, below the range, where her text says within. Verdict: supported.
\end{description}
\textbf{Conclusion.} Explained; both sides, high confidence. Our side's part was reading ``similar'' strictly only; \href{https://github.com/alectionary/vibe-math/blob/master/formalization/chen-rdd-mediation/simulation/statement_checks.py}{\nolinkurl{simulation/statement_checks.py}} reports every statement against her own figures, and statements 30 and 44 are reported here under both readings.

\subsection{A TNDE power gap at one cell set}\label{dis:tnde-power-gap}

\textbf{Her result.} TNDE power at cutoff 1, \cp{} $= .39$, $a = .59$, $h = 0$, $n = 500$ is .736 in her Figure 4.26. \textbf{Ours.} .698 in the same cells, 3.7 combined standard errors lower.
\begin{description}[leftmargin=1em,style=nextline]
\item[Our TNDE estimates are more variable (ours).] Prediction: our standard deviation would exceed her code's. Test: ours against the replay of her printed code on the 4,000 datasets it generates. Result: ratio 1.014, which lowers normal-approximation power by about .01, a quarter of the .039 gap. Verdict: ruled out.
\item[We misread her point (ours).] Prediction: neighbouring series in the same panel would read close to ours. Result: her three directly read $h$ series at $n = 500$ (the fourth is hidden) read .736-.762, all above ours. Verdict: ruled out.
\item[Monte Carlo noise in one prespecified comparison (neither).] Prediction: within about 2 combined standard errors. Result: 3.7, too large for one prespecified comparison; but this cell set was singled out after the fact as the largest gap among thousands of points, so the single-comparison criterion does not apply, and the next theory tests the chance explanation properly. Verdict: inconclusive.
\item[Her particular random draws (neither).] Her seeded datasets in these $h = 0$ cells give higher power than ours, so the gap is between two sets of random draws, not between procedures. Prediction: on the datasets her printed code generates, power computed the same way would exceed ours by most of the gap, and the gap would be no larger than expected for the most extreme of the thousands of points compared. Test: normal-approximation power from the mean and standard deviation of the estimates on those 4,000 datasets and on our 4,000. Result: .715 on hers against .687 on ours, a difference of .029 of the .039 gap (her .7362 against our .6975); mean TNDE bias is $-$.0004 in hers and $-$.0087 in ours; the remaining .010 is about 1.4 times her reading error of this point (.0071) and below one combined Monte Carlo and reading error (about .012), and a 3.7 standard-error gap is about what the most extreme of the 3,315 directly read rate points compared would show by chance. Verdict: supported.
\item[Her bootstrap differs in a detail from Appendix B (source).] Prediction: her intervals would be narrower in these cells. Untested: her bootstrap is unseeded and her intervals are not reported, so the prediction cannot be checked; the gap is one cell set and changes no conclusion.
\end{description}
\textbf{Conclusion.} Explained as Monte Carlo variation between her random draws and ours, singled out as the most extreme of many comparisons; neither side's error, medium confidence.

\subsection{The 5,000-case recommendation overgeneralizes}\label{dis:n5000-overgeneralized}

\textbf{Her statement.} $n = 5{,}000$ is the tentative minimum for accurate Type I error and power (Chapter 5). \textbf{Ours.} At $n = 5{,}000$, 86.3\% of zero-effect cells have Type I error in Bradley's range; effects whose two paths are both zero stay conservative at every $n$.
\begin{description}[leftmargin=1em,style=nextline]
\item[Our Type I error is wrong when both paths are zero (ours).] Prediction: her own figure would show these cells within range at $n = 5{,}000$. Test: PNIE Type I error with $a = b = 0$ by $n$, ours and her Figure 4.17. Result: hers .0014-.0020 and ours .0015-.0020 at every $n$. Verdict: ruled out.
\item[The recommendation does not account for both-zero cells (source).] Prediction: her own figure would show them below .025 at $n = 5{,}000$ and 10,000. Result: .0015 and .0014 in her Figure 4.17. Verdict: supported.
\end{description}
\textbf{Conclusion.} Explained; source-side, high confidence.

\subsection{A Table 4.3 \texorpdfstring{$t$}{t} value}\label{dis:table-4-3-t-value}

\textbf{Her table.} Table 4.3, cutoff 3, child baseline age: coefficient $-0.127$, standard error 0.061, $t = -2.252$. \textbf{Ours.} $-0.127/0.061 = -2.082$ (2.057-2.107 in absolute value within the printed rounding); the other eight rows agree.
\begin{description}[leftmargin=1em,style=nextline]
\item[We transcribed the row wrongly (ours).] Test: re-read from the PDF text layer. Result: $-0.127$, 0.061, $-2.252$. Verdict: ruled out.
\item[Rounding (neither).] Test: bounds of $|$coefficient$|$/SE over the rounding intervals. Result: 2.057-2.107; 2.252 is outside. Verdict: ruled out.
\item[A printed typo (source).] Test: solve for the number the printed $t$ implies. Result: a coefficient of $-0.137$ (one digit) gives $t = -2.25$; which number is wrong cannot be told, and the row is significant at .05 either way. Verdict: supported.
\end{description}
\textbf{Conclusion.} Explained; source-side, high confidence; no conclusion changes.

\subsection{Opposite significance conclusions in the Head Start analysis}\label{dis:significance-conclusions}

\textbf{Her statement.} On the Head Start data, 95\% intervals for the RDD estimates ``made opposite significance conclusions than the RCT estimates'' (abstract); ``if the RCT was significant, the RDD estimate would not be significant'' (printed p.~143). Her printed p.~81 also says that PNDE and TNDE ``were not significant''. \textbf{Ours.} From her own printed Table 4.7 intervals, 10 of the 15 trial-versus-RDD comparisons reach opposite conclusions and 5 the same (PNDE and TNDE at cutoff 3, MINT at all three cutoffs). Where the trial estimate is significant (PNIE, TNIE) the RDD estimate is not, in all 6 cells; but in 4 cells (PNDE and TNDE at cutoffs 1-2) the RDD interval excludes zero where the trial interval does not.
\begin{description}[leftmargin=1em,style=nextline]
\item[We transcribed the intervals wrongly (ours).] Prediction: the PDF text layer and the rendered page would show different bounds. Test: re-read Table 4.7 from the text layer (PDF pp. 97-98) and the rendered page. Result: the bounds match both; the four RDD direct-effect intervals that exclude zero read [-1.592, -0.413], [-1.696, -0.375], [-1.606, -0.389] and [-1.689, -0.396]. Verdict: ruled out.
\item[Her summary generalizes from the indirect effects (source).] Her summary may extend the indirect-effect pattern to all effects, with the p.~81 sentence referring to the trial's direct effects only. Prediction: the pattern would hold in every cell where the trial estimate is significant and fail elsewhere, and the trial PNDE and TNDE intervals would include zero. Test: significance of every printed interval. Result: it holds in all 6 cells where the trial estimate is significant and fails in the 9 where it is not (4 opposite in the other direction, 5 the same); the trial PNDE and TNDE intervals include zero, while the RDD ones at cutoffs 1-2 do not, so the p.~81 sentence is true only of the trial. Her own Chapter 4 reading of the table (printed p.~85) agrees with ours cell by cell, calling the PNIE and TNIE RDD intervals Type II errors and the PNDE and TNDE ones at cutoffs 1-2 Type I errors. Verdict: supported.
\end{description}
\textbf{Conclusion.} Explained; source-side, high confidence. Her statement is right for the indirect effects and overgeneralized for the others; the overgeneralization is in her abstract and p.~143 summary only, not in her Chapter 4 reading of the table.

\section{Assumptions we supplied}\label{sec:assumptions}

Each choice below was ours, not stated by Chen.

\subsection{Population values from the effect definitions}\label{sec:ass-pop}
Chen does not print the population values she compared estimates with. We computed them from her effect definitions and her data-generating model, giving PNDE $= c'$ because the mediator intercept is 0 (Section~\ref{dis:pnde-definition}). This choice affects bias and coverage, not rejection rates.

\subsection{Numeric readings of her words}\label{sec:ass-words}
To test her prose we fixed thresholds: for example, ``near zero'' Type I error $<$ .025, ``close to zero'' power $<$ .10, and ``around .80'' between .70 and .90. A second reading moves every threshold .01 in her favor; the slack column of \href{https://github.com/alectionary/vibe-math/blob/master/formalization/chen-rdd-mediation/results/statement_checks.csv}{\nolinkurl{results/statement_checks.csv}} shows which verdicts depend on these choices.

\subsection{Averaging over unnamed factors}\label{sec:ass-avg}
A statement or figure that does not name a factor is compared with our rates averaged over that factor, as her figures appear to do. Her aggregation code is not printed.

\subsection{Scalar summaries of her tables}\label{sec:ass-contrasts}
The headline contrasts E1-E3 average her printed Table 4.8 and 4.9 rows (six direct-effect rows; fifteen rows), and E4 averages one design cell over $b$, \cp{} and cutoffs. These summaries are ours; the rows themselves are in \href{https://github.com/alectionary/vibe-math/blob/master/formalization/chen-rdd-mediation/source-values/table_4_8.csv}{\nolinkurl{source-values/table_4_8.csv}} and \href{https://github.com/alectionary/vibe-math/blob/master/formalization/chen-rdd-mediation/source-values/table_4_9.csv}{\nolinkurl{source-values/table_4_9.csv}}.

\subsection{The tolerance for her figure value}\label{sec:ass-tol}
For E4 the tolerance is a measurement tolerance: the mean reading error of her two directly read Figure 4.28 points (cutoffs 2-3), 0.0091, rounded down, and not set in the plan. It is used instead of a rounding bound because the value is not printed as a number. Her cutoff-1 series is hidden under the others in that cell and is not read, so the contrast uses cutoffs 2-3 on both sides.

\section{Ambiguities that were resolved}\label{sec:ambiguities}

\subsection{Which side of the cutoff is treated}\label{sec:amb-cutoff}
\textbf{Ambiguity.} Her Chapter 3, Table 3.1 and the note to Table 4.8 call cutoff index 1 ``25\% treatment and 75\% control'', but her Appendix A code assigns treatment at or above the 25th percentile, so index 1 has 75\% treated. \textbf{Check.} Her code and her Head Start Table 4.2. \textbf{Finding.} The code treats the upper side; her Head Start analysis treats the lower side (treated shares .356, .717, .856), as her text says. \textbf{Implication.} We report the simulation by the code, because her printed code defines the datasets she simulated; Section~\ref{dis:population-values} tests whether the reversal entered her coverage step (inconclusive). Evidence: \href{https://github.com/alectionary/vibe-math/blob/master/formalization/chen-rdd-mediation/source-values/table_3_1.csv}{\nolinkurl{source-values/table_3_1.csv}}, \href{https://github.com/alectionary/vibe-math/blob/master/formalization/chen-rdd-mediation/source-values/table_4_2.csv}{\nolinkurl{source-values/table_4_2.csv}}, \href{https://github.com/alectionary/vibe-math/blob/master/formalization/chen-rdd-mediation/source-values/README.md}{\nolinkurl{source-values/README.md}}.

\subsection{The mediator intercept in PNDE}\label{sec:amb-pnde}
\textbf{Ambiguity.} Her text gives PNDE as \cp{} + $h$; the general definition is $c' + h\,i_{0M}$. \textbf{Check.} The replay's fitted intercept, and the identity on her Head Start estimates. \textbf{Finding.} The intercept is 0 in her simulation, and her real-data code used the fitted intercept. \textbf{Implication.} The population PNDE is \cp{} (Section~\ref{dis:pnde-definition}).

\subsection{The standardized-bias divisor}\label{sec:amb-sb}
\textbf{Ambiguity.} Equation 3.12 and her tables disagree on the divisor. \textbf{Check and finding.} All 24 printed values use the estimate's standard error. \textbf{Implication.} S5 is judged with that divisor (Section~\ref{dis:headstart-sb-divisor}; \href{https://github.com/alectionary/vibe-math/blob/master/formalization/chen-rdd-mediation/secondary/table_arithmetic.py}{\nolinkurl{secondary/table_arithmetic.py}}).

\subsection{``Similar patterns'' across cutoffs}\label{sec:amb-similar}
\textbf{Ambiguity.} Statements 30 and 44 say later cutoffs show ``similar'' patterns. \textbf{Check and finding.} Under the strict reading both her figures and our run fail; under a same-shape reading both pass. \textbf{Implication.} Both readings are reported (Section~\ref{dis:prose-versus-own-figure}).

\section{What remains unresolved}\label{sec:unresolved}

\textbf{The per-case population-value rule (partly explained).} Candidates: one wrong value per condition; values right for control and wrong for treated cases; a step written for her text's reversed treatment labels; noise. Attempts: the first and the last were ruled out by her coverage plateau and bias fits; the per-case rule reproduces cutoffs 1-2 but gives $\lambda = 1.76$ rather than about 1.9-2.0 at cutoff 3; the label reversal predicts the same plateau and cannot be separated from it. Ranking of the remaining causes: a per-case rule with a cutoff-specific detail we have not identified, then the label reversal. The artifact that would settle the largest one is her aggregation code or result files with her population values, which are not public.

\textbf{The mechanism that removed the interaction.} The two candidates differ only in the $y$ intercept; her $y$ intercepts or result files would separate them.

\textbf{Our small positive imbalance.} Our run puts the true value above the interval slightly more often than below it (0.33 points on average); we did not investigate why.

\textbf{The TNDE power gap's last part.} About .010 of the .039 gap is not attributed to her random draws; it is about 1.4 times her reading error of that point (.0071) and below one combined Monte Carlo and reading error (about .012), and the bootstrap-detail theory cannot be tested without her intervals.

\textbf{RQ2 and the announced power analysis.} S4 and S6 are blocked by data access; S8 was announced but never reported, so it is not established and nothing blocks it.

\section{Limitations}\label{sec:limitations}

\begin{enumerate}[leftmargin=1.5em]
\item\label{lim:0} The bootstrap-based results (coverage, Type I error, power) come from the Python run with our own seeds; Chen's bootstrap is unseeded, so they agree with hers on summaries only. The R replay reproduces exactly the datasets and point estimates her printed, seeded code generates (all 1,000 replications per condition); her published results evidently came from data without the interaction wherever $h > 0$ (an inference from her figures; item 2), which the printed code does not produce, so the replay is not her published runs (H1-H4, S1-S3).
\item\label{lim:1} Chen's aggregation code, population values and result files are not printed. The mechanisms behind her bias, coverage and MINT results are inferred from her digitized figures and printed tables; they reproduce her patterns but are not read from her code (H1, H2, H3, S1).
\item\label{lim:2} Figures 4.17-4.38 were read from images, with reading errors up to about .02 on Type I error and power, median .03 and up to about .05 on coverage (Figures 4.31 and 4.35); points hidden behind another series are not readings: the digitizer fills them in from the visible series (status occluded), and they are not used for statement verdicts, figure-agreement rates or E4 (H3, S2, S3).
\item\label{lim:3} RQ2 (HSIS) cannot be rerun: the restricted-use data are not available to this project (S4, S6); only the printed arithmetic and identities were checked (S5).
\item\label{lim:4} The announced power analysis by the standard deviation of $Y$ is not reported, so it cannot be assessed (S8).
\item\label{lim:5} The dissertation labels the cutoffs inconsistently: the code's index 1 is the 25th percentile (75\% treated), while Chapter 3, the Table 4.8 note and her Head Start analysis have treatment below the cutoff; we report by the code (H1, H2).
\item\label{lim:6} Chen suggests global estimation with an incorrect functional form as a plausible explanation of the bias \citep[printed p.~147]{chen2021}; in her data-generating model the assignment variable does not enter $M$ or $Y$, so the fitted linear models are correctly specified (H1).
\end{enumerate}
Claim-specific limits are stated with each claim: the H3 value is read from an image at cutoffs 2-3 only, with mean reading error 0.0091; Chen calls the H4 recommendation tentative while we judge it literally; the zero-estimate pattern of S1 does not fit every printed entry; and the S5 divisor discrepancy is between her text and her tables, with her intent unknown. S4, S6, S7 and S8 are unassessed.

\section{Code and source availability}\label{sec:availability}

The source is \citet{chen2021}, from the University of Nebraska-Lincoln Digital Commons record \url{https://digitalcommons.unl.edu/cehsdiss/403/} (SHA-256 \texttt{9206b3d8\ldots e557}, 6,720,692 bytes, 203 pages). As a document by another author, it is kept out of the repository; the values read from it are committed under \href{https://github.com/alectionary/vibe-math/tree/master/formalization/chen-rdd-mediation/source-values}{\nolinkurl{source-values/}} and in the digitized result files, and its figures and tables are inventoried in \href{https://github.com/alectionary/vibe-math/blob/master/formalization/chen-rdd-mediation/source-inventory.json}{\nolinkurl{source-inventory.json}}. Her simulation code is printed in her Appendices A-B and replayed by \href{https://github.com/alectionary/vibe-math/blob/master/formalization/chen-rdd-mediation/r-reproduction/chen_exact.R}{\nolinkurl{r-reproduction/chen_exact.R}} (\href{https://github.com/alectionary/vibe-math/blob/master/formalization/chen-rdd-mediation/r-reproduction/README.md}{\nolinkurl{r-reproduction/README.md}}). The source-data availability ledger for RQ2, checked on 2026-11-08 (UTC) and kept in \href{https://github.com/alectionary/vibe-math/blob/master/formalization/chen-rdd-mediation/notes/availability.md}{\nolinkurl{notes/availability.md}}, is:

{\footnotesize
\begin{longtable}{@{}>{\raggedright\arraybackslash}p{0.33\textwidth}>{\raggedright\arraybackslash}p{0.1\textwidth}>{\raggedright\arraybackslash}p{0.5\textwidth}@{}}
\toprule
Place checked & Date & Result\\
\midrule
Chen (2021), Chapter 3, Dataset & 2026-11-08 & Uses the HSIS restricted-use dataset from ICPSR, accessed under agreement; 3,721 children after excluding crossovers and no-shows, 2,697 with a wave-1 spelling score.\\
\href{https://www.icpsr.umich.edu/web/ICPSR/studies/29462}{ICPSR 29462}, Head Start Impact Study 2002-2006 & 2026-11-08 & The study page offers a public download after free sign-in and lists restricted-use files that need a signed Restricted Data Use Agreement; its wording suggests public-use data files, but the next row finds none.\\
ICPSR 29462 public download, version V7 (2018-02-08) & 2026-11-08 & 58 files, 454 MB unpacked: user guide, overview, codebooks and questionnaires for DS0001-DS0030. No data files.\\
\href{https://www.icpsr.umich.edu/web/ICPSR/studies/36968}{ICPSR 36968}, HSIS with Center Analysis File & 2026-11-08 & Related restricted collection; not checked further.\\
Restricted-use application by the repository owner & 2026-11-08 & Applied for the HSIS restricted-use data and could not obtain it (recorded in \href{https://github.com/alectionary/vibe-math/blob/master/formalization/chen-rdd-mediation/REPLICATION-PLAN.yaml}{\nolinkurl{REPLICATION-PLAN.yaml}}).\\
Chen (2021), appendices & 2026-11-08 & No analysis data or input for RQ2; Appendices A-B cover only the simulation.\\
\bottomrule
\end{longtable}}

Her printed summaries do not substitute for the data: Table 4.2 gives within-group means and standard deviations and Tables 4.6-4.7 give estimates with standard errors, but no covariances of treatment with mediator or outcome are printed, so the paths cannot be refitted (\href{https://github.com/alectionary/vibe-math/blob/master/formalization/chen-rdd-mediation/REPLICATION-PLAN.yaml}{\nolinkurl{REPLICATION-PLAN.yaml}}). The project owner applied for the restricted-use data and could not obtain them (ledger row ``Restricted-use application'').

\section{Next steps}\label{sec:next}

The artifacts named in Section~\ref{sec:unresolved} would settle the open questions: her aggregation code or result files would show the per-case population-value rule and, with her $y$ intercepts, which mechanism removed the interaction; the check would compare her stored population values and intervals with our cells cell by cell. Restricted-use HSIS access under an ICPSR agreement would allow RQ2 (S4, S6) to be rerun with her printed model. Within our own evidence, the next check is the cause of our 0.33-point imbalance, for example by comparing percentile with bias-corrected intervals on a subset of cells.

\bibliographystyle{plainnat}
\bibliography{references}

\appendix

\section{Artifacts and reproduction}\label{app:artifacts}

All commands run from the study folder \href{https://github.com/alectionary/vibe-math/tree/master/formalization/chen-rdd-mediation}{\nolinkurl{formalization/chen-rdd-mediation}} with Python 3 and NumPy (scikit-learn for the estimator check). The code, committed outputs and evidence cited here are those on the \texttt{master} branch, where every link points. The source PDF must be downloaded from the record above into \texttt{source/chen2021.pdf} for the digitizing and packet steps. The per-condition simulation outputs (\texttt{simulation/work/}) are not committed; \href{https://github.com/alectionary/vibe-math/blob/master/formalization/chen-rdd-mediation/README.md}{\nolinkurl{README.md}} (section ``Reproduce'') gives the full recipe.

{\footnotesize\def\UrlBreaks{\do\/\do\_\do\-}\def\UrlBigBreaks{}
\begin{longtable}{@{}>{\raggedright\arraybackslash}p{0.2\textwidth}>{\raggedright\arraybackslash}p{0.33\textwidth}>{\raggedright\arraybackslash}p{0.41\textwidth}@{}}
\toprule
Claim or result & File & Command (output)\\
\midrule\endhead
Estimator check & \href{https://github.com/alectionary/vibe-math/blob/master/formalization/chen-rdd-mediation/simulation/validate_estimator.py}{\nolinkurl{simulation/validate_estimator.py}} & \texttt{\detokenize{cd simulation && python3 validate_estimator.py}} (prints \texttt{\detokenize{PASS_ESTIMATOR_VALIDATION}})\\
H1-H4, S1-S3 simulation & \href{https://github.com/alectionary/vibe-math/blob/master/formalization/chen-rdd-mediation/simulation/chen_sim.py}{\nolinkurl{simulation/chen_sim.py}} & \texttt{\detokenize{cd simulation && OPENBLAS_NUM_THREADS=1 python3 chen_sim.py run --out work --workers 2}} (1,280 files; about 13.7 core-hours; resumes)\\
H1, H2, S1 tables & \href{https://github.com/alectionary/vibe-math/blob/master/formalization/chen-rdd-mediation/results/table48_ours.csv}{\nolinkurl{results/table48_ours.csv}}, \href{https://github.com/alectionary/vibe-math/blob/master/formalization/chen-rdd-mediation/results/table49_ours.csv}{\nolinkurl{results/table49_ours.csv}}, \href{https://github.com/alectionary/vibe-math/blob/master/formalization/chen-rdd-mediation/results/cells.csv}{\nolinkurl{results/cells.csv}} & \texttt{\detokenize{cd simulation && python3 aggregate.py work ../results}} (\href{https://github.com/alectionary/vibe-math/blob/master/formalization/chen-rdd-mediation/simulation/aggregate.py}{\nolinkurl{simulation/aggregate.py}}; also type1, power, summary)\\
H3 MINT power & \href{https://github.com/alectionary/vibe-math/blob/master/formalization/chen-rdd-mediation/results/summary.json}{\nolinkurl{results/summary.json}}, \href{https://github.com/alectionary/vibe-math/blob/master/formalization/chen-rdd-mediation/results/power.csv}{\nolinkurl{results/power.csv}} & as above\\
H4, S2, S3 statements & \href{https://github.com/alectionary/vibe-math/blob/master/formalization/chen-rdd-mediation/results/statement_checks.csv}{\nolinkurl{results/statement_checks.csv}}, \href{https://github.com/alectionary/vibe-math/blob/master/formalization/chen-rdd-mediation/results/type1.csv}{\nolinkurl{results/type1.csv}} & \texttt{\detokenize{cd simulation && python3 statement_checks.py ../results/cells.csv > ../results/statement_checks.csv}}\\
Exact replay of her code & \href{https://github.com/alectionary/vibe-math/blob/master/formalization/chen-rdd-mediation/r-reproduction/chen_exact.R}{\nolinkurl{r-reproduction/chen_exact.R}}, \href{https://github.com/alectionary/vibe-math/blob/master/formalization/chen-rdd-mediation/r-reproduction/out-full/condition_summary.csv}{\nolinkurl{r-reproduction/out-full/condition_summary.csv}} & \texttt{\detokenize{cd r-reproduction && Rscript chen_exact.R 1000 out-full summary}} (R 4.1.3, 2021-08-02 packages; about 3.4 hours)\\
Replay against our run & \href{https://github.com/alectionary/vibe-math/blob/master/formalization/chen-rdd-mediation/results/r_full_comparison.json}{\nolinkurl{results/r_full_comparison.json}} & \texttt{\detokenize{python3 secondary/r_full_comparison.py}} (\href{https://github.com/alectionary/vibe-math/blob/master/formalization/chen-rdd-mediation/secondary/r_full_comparison.py}{\nolinkurl{secondary/r_full_comparison.py}})\\
Digitized bias & \href{https://github.com/alectionary/vibe-math/blob/master/formalization/chen-rdd-mediation/results/digitized_bias.csv}{\nolinkurl{results/digitized_bias.csv}}, \href{https://github.com/alectionary/vibe-math/blob/master/formalization/chen-rdd-mediation/results/bias_structure.csv}{\nolinkurl{results/bias_structure.csv}} & \texttt{\detokenize{python3 secondary/digitize_bias_figures.py source/chen2021.pdf}}; \href{https://github.com/alectionary/vibe-math/blob/master/formalization/chen-rdd-mediation/secondary/bias_structure.py}{\nolinkurl{secondary/bias_structure.py}}\\
Digitized rates and imbalance & \href{https://github.com/alectionary/vibe-math/blob/master/formalization/chen-rdd-mediation/results/digitized_rates.csv}{\nolinkurl{results/digitized_rates.csv}}, \href{https://github.com/alectionary/vibe-math/blob/master/formalization/chen-rdd-mediation/results/digitized_imbalance.csv}{\nolinkurl{results/digitized_imbalance.csv}}, \href{https://github.com/alectionary/vibe-math/blob/master/formalization/chen-rdd-mediation/results/digitize_crosscheck.csv}{\nolinkurl{results/digitize_crosscheck.csv}}, \href{https://github.com/alectionary/vibe-math/blob/master/formalization/chen-rdd-mediation/results/rates_agreement.csv}{\nolinkurl{results/rates_agreement.csv}} & \texttt{\detokenize{python3 secondary/digitize_raster_figures.py source/chen2021.pdf}}; \href{https://github.com/alectionary/vibe-math/blob/master/formalization/chen-rdd-mediation/secondary/digitize_imbalance_figures.py}{\nolinkurl{secondary/digitize_imbalance_figures.py}}; \href{https://github.com/alectionary/vibe-math/blob/master/formalization/chen-rdd-mediation/secondary/digitize_crosscheck.py}{\nolinkurl{secondary/digitize_crosscheck.py}}; \href{https://github.com/alectionary/vibe-math/blob/master/formalization/chen-rdd-mediation/secondary/rates_agreement.py}{\nolinkurl{secondary/rates_agreement.py}}\\
Disagreement tests & \href{https://github.com/alectionary/vibe-math/blob/master/formalization/chen-rdd-mediation/results/disagreement_tests.json}{\nolinkurl{results/disagreement_tests.json}} & \texttt{\detokenize{python3 secondary/disagreement_tests.py}} (\href{https://github.com/alectionary/vibe-math/blob/master/formalization/chen-rdd-mediation/secondary/disagreement_tests.py}{\nolinkurl{secondary/disagreement_tests.py}})\\
S1 zero-estimate pattern & \href{https://github.com/alectionary/vibe-math/blob/master/formalization/chen-rdd-mediation/secondary/zero_estimate_check.csv}{\nolinkurl{secondary/zero_estimate_check.csv}} & \texttt{\detokenize{python3 secondary/zero_estimate_check.py > secondary/zero_estimate_check.csv}}\\
S5 arithmetic & \href{https://github.com/alectionary/vibe-math/blob/master/formalization/chen-rdd-mediation/secondary/table_arithmetic.csv}{\nolinkurl{secondary/table_arithmetic.csv}} & \texttt{\detokenize{cd secondary && python3 table_arithmetic.py}}\\
S9 significance conclusions & \href{https://github.com/alectionary/vibe-math/blob/master/formalization/chen-rdd-mediation/results/significance_conclusions.json}{\nolinkurl{results/significance_conclusions.json}} & \texttt{\detokenize{python3 secondary/significance_conclusions.py}} (\href{https://github.com/alectionary/vibe-math/blob/master/formalization/chen-rdd-mediation/secondary/significance_conclusions.py}{\nolinkurl{secondary/significance_conclusions.py}}; input \href{https://github.com/alectionary/vibe-math/blob/master/formalization/chen-rdd-mediation/source-values/table_4_7.csv}{\nolinkurl{source-values/table_4_7.csv}})\\
S5 identities & \href{https://github.com/alectionary/vibe-math/blob/master/formalization/chen-rdd-mediation/results/empirical_consistency.json}{\nolinkurl{results/empirical_consistency.json}}, \href{https://github.com/alectionary/vibe-math/blob/master/formalization/chen-rdd-mediation/results/table_identities.json}{\nolinkurl{results/table_identities.json}} & \texttt{\detokenize{python3 secondary/empirical_consistency.py}} (\href{https://github.com/alectionary/vibe-math/blob/master/formalization/chen-rdd-mediation/secondary/empirical_consistency.py}{\nolinkurl{secondary/empirical_consistency.py}})\\
Evidence packet & \href{https://github.com/alectionary/vibe-math/blob/master/formalization/chen-rdd-mediation/evidence-packet.json}{\nolinkurl{evidence-packet.json}} & \texttt{\detokenize{python3 build_packet.py}} (\href{https://github.com/alectionary/vibe-math/blob/master/formalization/chen-rdd-mediation/build_packet.py}{\nolinkurl{build_packet.py}}; needs the source PDF and pdftotext)\\
This paper & \href{https://github.com/alectionary/vibe-math/blob/master/formalization/chen-rdd-mediation/paper/mediation_in_regression_discontinuity__a_replication_of_chen_2021.tex}{\nolinkurl{paper/mediation_in_regression_discontinuity__a_replication_of_chen_2021.tex}} & \texttt{\detokenize{cd paper && pdflatex; bibtex; pdflatex; pdflatex}} (writer record \href{https://github.com/alectionary/vibe-math/blob/master/formalization/chen-rdd-mediation/writer-manifest.json}{\nolinkurl{writer-manifest.json}})\\
\bottomrule
\end{longtable}}

\Needspace{12\baselineskip}
\section{Compliance ledger}\label{app:ledger}

{\footnotesize
\begin{longtable}{@{}>{\raggedright\arraybackslash}p{0.28\textwidth}>{\raggedright\arraybackslash}p{0.1\textwidth}>{\raggedright\arraybackslash}p{0.52\textwidth}@{}}
\toprule
Requirement & Status & Evidence\\
\midrule\endhead
Engine (preflight) & departed & The source's software is R; the bootstrap-based results come from the recorded Python fallback, which the preflight accepts with its independent validation and the owner's decision of 2026-11-09 (\href{https://github.com/alectionary/vibe-math/blob/master/formalization/chen-rdd-mediation/simulation/fallback.json}{\nolinkurl{simulation/fallback.json}}). Her printed R code was replayed exactly for the datasets and point estimates it generates (Section~\ref{sec:engines}). Justification, as in Section~\ref{sec:engines}: her bootstrap is unseeded, so no engine can match it draw for draw, and a full R rerun of the bootstrap (about 4 billion \texttt{lm()} fits) was declined.\\
Source software & departed & Python was first used without asking the owner; the owner then agreed. Justification as above.\\
Full design run & followed & 3,840 of 3,840 cells at 1,000 replications and 500 resamples (\href{https://github.com/alectionary/vibe-math/blob/master/formalization/chen-rdd-mediation/notes/run-record.md}{\nolinkurl{notes/run-record.md}}).\\
Verdict first, limits next & followed & Abstract and Section~\ref{sec:intro}.\\
Page-1 headline table with magnitude and direction & followed & Table~1.\\
Every disagreement with theories and tests & followed & Section~\ref{sec:disagreements}; \href{https://github.com/alectionary/vibe-math/blob/master/formalization/chen-rdd-mediation/DISAGREEMENTS.yaml}{\nolinkurl{DISAGREEMENTS.yaml}}.\\
Figures digitized with reading error & followed & Section~\ref{sec:digitize}.\\
Blocked parts not called nonreplication & followed & Section~\ref{sec:rq2}.\\
Software cited with versions & departed & R, NumPy and scikit-learn cited with versions; PyMuPDF and pandas versions used for digitizing were not recorded.\\
Citations resolve; every reference used & followed & \texttt{\detokenize{check_citations.py}} on this paper and \texttt{paper/build.log}.\\
Code and data links point at master & followed & Appendix~\ref{app:artifacts}.\\
No authorship statement by a language model & followed & None in the text.\\
\bottomrule
\end{longtable}}

\section{Claim and section inventory}\label{app:inventory}

Every claim the evidence packet records, with its status and where it is discussed, and every section of the dissertation with the claims it carries.

{\footnotesize
\begin{longtable}{@{}>{\raggedright\arraybackslash}p{0.06\textwidth}>{\raggedright\arraybackslash}p{0.42\textwidth}>{\raggedright\arraybackslash}p{0.16\textwidth}>{\raggedright\arraybackslash}p{0.26\textwidth}@{}}
\toprule
Claim & Statement \citep{chen2021} & Status & Discussed\\
\midrule\endhead
H1 & Direct-effect estimates (PNDE, TNDE) are biased, mostly negatively (Table 4.8). & not replicated & Table 1; Sections~\ref{sec:h1}, \ref{dis:h-absent}-\ref{dis:pnde-definition}\\
H2 & Percentile bootstrap intervals cover well below .95, true value above the interval (Table 4.9). & not replicated & Table 1; Sections~\ref{sec:h2}, \ref{dis:population-values}\\
H3 & MINT power is near zero at every sample size (Figure 4.28). & not replicated & Table 1; Sections~\ref{sec:h3}, \ref{dis:h-absent}\\
H4 & Tentatively, $n = 5{,}000$ gives accurate Type I error and power above .80 at any cutoff. & partly replicated & Table 1; Sections~\ref{sec:h4}, \ref{dis:n5000-overgeneralized}\\
S1 & Indirect and MINT bias within [-.10, .10] in 51.88-87.50\% of conditions. & not replicated & Section~\ref{sec:s1}\\
S2 & Type I error statements for all five effects (Figures 4.17-4.21). & partly replicated & Section~\ref{sec:s2}\\
S3 & Power statements for all five effects (Figures 4.22-4.28). & partly replicated & Section~\ref{sec:s3}\\
S4 & HSIS synthetic-RDD estimates typically smaller than RCT estimates. & blocked & Section~\ref{sec:rq2}\\
S5 & Tables 4.6-4.7 differences and standardized biases. & replicated & Sections~\ref{sec:rq2}, \ref{dis:headstart-sb-divisor}\\
S6 & HSIS assumption checks show slight violations. & blocked & Section~\ref{sec:rq2}\\
S7 & Integrated assumption table answers RQ1 (Table 4.1). & not established & Section~\ref{sec:s7s8}\\
S8 & Additional power analysis by the SD of $Y$. & not established: announced, never reported & Section~\ref{sec:s7s8}\\
S9 & HSIS RDD intervals give opposite significance conclusions to the RCT. & partly replicated & Sections~\ref{sec:rq2}, \ref{dis:significance-conclusions}\\
\bottomrule
\end{longtable}}

{\footnotesize
\begin{longtable}{@{}>{\raggedright\arraybackslash}p{0.4\textwidth}>{\raggedright\arraybackslash}p{0.15\textwidth}>{\raggedright\arraybackslash}p{0.39\textwidth}@{}}
\toprule
Dissertation section & Claims & Note\\
\midrule\endhead
Front matter & none & The abstract restates the Chapter 4-5 findings.\\
Chapter 1, Introduction & H1-H4, S7 & Research questions and design.\\
Chapter 2, Literature review & none & Effect definitions used for the population values.\\
Chapter 3, Methods and procedures & S5, S8 & Design, equation 3.12, announced power analysis.\\
Chapter 4, Results & H1-H3, S1-S7, S9 & RQ1 Table 4.1; RQ2 Tables 4.2-4.7; RQ3 Tables 4.8-4.9, Figures 4.12-4.38.\\
Chapter 5, Discussion & H1-H4, S4, S6, S9 & Conclusions.\\
References & none & Bibliography only.\\
Appendix A, data generation (printed pp. 175-178) & H1, H2 & Replayed exactly.\\
Appendix B, estimation and bootstrap (printed pp. 179-189) & H1-H3 & Point estimates replayed; bootstrap unseeded.\\
\bottomrule
\end{longtable}}

\end{document}
